\ChapterImageStar[cap:cumplimiento-objetivos]{Cumplimiento de objetivos}{./images/fondo.png}\label{cap:cumplimiento-objetivos}
\mbox{}\\

El presente capítulo documenta el cumplimiento sistemático de los objetivos establecidos para la investigación, evidenciando cómo cada una 
de las fases desarrolladas contribuyó al logro del objetivo general orientado a aspectos de seguridad para la infraestructura de 
computación en la nube del Grupo de Investigación en Redes, Información y Distribución (\GRID) de la Universidad del Quindío. En este
apartado se presenta una evaluación del grado de cumplimiento de los objetivos específicos, respaldada por los resultados obtenidos 
durante las etapas de caracterización, estudio de mapeo sistemático, análisis DAR, selección tecnológica y modelado de la arquitectura de 
defensa en profundidad propuesta.

\section{Cumplimiento del objetivo general}
\noindent

\textbf{Objetivo general}: Especificar una arquitectura de seguridad para la infraestructura de la computación en la
nube del grupo de investigación \GRID.

\textbf{Estado de cumplimiento}: \textbf{CUMPLIDO COMPLETAMENTE}

El objetivo general se cumplió satisfactoriamente mediante la especificación de una arquitectura de defensa en profundidad orientada al 
fortalecimiento de la supervisión, monitoreo y protección de la infraestructura cloud del grupo \GRID. La investigación permitió 
identificar las necesidades, problemas y oportunidades relacionadas con la seguridad de la infraestructura, así como seleccionar y modelar 
un conjunto de tecnologías interoperables capaces de proporcionar capacidades de correlación de eventos, monitoreo continuo, detección de 
intrusiones y respuesta ante incidentes.

Como resultado, se diseñó una arquitectura multicapa soportada en tecnologías como SonicWall NSA 6600, Suricata y Wazuh, integradas 
mediante mecanismos de segmentación, análisis de tráfico y centralización de registros de seguridad. Asimismo, el desarrollo del estudio de 
mapeo sistemático, el análisis DAR y el modelado arquitectónico en Archimate y Draw.io permitieron consolidar una propuesta alineada con 
controles de las normas \ISO/IEC 27001 e \ISO/IEC 27017, proporcionando una arquitectura reproducible, escalable y adaptable a entornos 
académicos con necesidades similares de protección y monitoreo de infraestructura.

\textbf{Evidencia del cumplimiento}:
\begin{itemize}
	\item Implementación funcional de las tecnologías Suricata y Wazuh dentro del entorno prototipo asociado a la infraestructura del \GRID.
	\item Definición de aspectos de seguridad alineados con controles de la \ISO/IEC 27001 y la \ISO/IEC 27017.
	\item Validación operativa de mecanismos de monitoreo continuo, correlación de eventos y respuesta activa ante incidentes de seguridad.
	\item Modelado arquitectónico completo de la arquitectura de defensa en profundidad mediante Archimate y Draw.io.
	\item Documentación técnica de la arquitectura propuesta, incluyendo fichas técnicas, segmentación de red y alineación con controles de la \ISO/IEC 27001 y la \ISO/IEC 27017.
\end{itemize}

\section{Cumplimiento de objetivos específicos}
\noindent

\subsection{Objetivo específico 1}
\noindent

\textbf{Enunciado}: Identificar necesidades, problemas y oportunidades relacionados con la ciberseguridad para la infraestructura de 
la computación en la nube del Grupo \GRID.

\textbf{Estado de cumplimiento}: \textbf{CUMPLIDO COMPLETAMENTE}

El objetivo general se cumplió satisfactoriamente mediante el análisis de las necesidades, problemas y oportunidades presentes en la 
infraestructura de computación en la nube del grupo \GRID, relacionadas con aspectos de ciberseguridad, monitoreo y gestión de controles de 
seguridad. A partir de la caracterización realizada y del análisis de la información obtenida, fue posible evidenciar la necesidad de 
alinear la gestión de seguridad con estándares internacionales como la \ISO/IEC 27001 y la \ISO/IEC 27017, así como centralizar la 
administración y supervisión de los controles de seguridad implementados dentro de la infraestructura.

Asimismo, se identificaron problemáticas asociadas a la baja adopción de normas y buenas prácticas de seguridad de la información, además 
de la dispersión y baja formalización de controles, lo cual genera brechas relacionadas con monitoreo, trazabilidad y protección de los 
activos de información del entorno \GRID. Estas condiciones evidenciaron la necesidad de fortalecer los mecanismos de supervisión, 
correlación de eventos y gestión centralizada de incidentes de seguridad.

De igual manera, el análisis permitió identificar oportunidades de mejora orientadas a la implementación de mecanismos de Gestión Unificada 
de Amenazas (UTM), la integración de controles de seguridad dentro de una arquitectura de defensa en profundidad y el fortalecimiento de la 
alineación con los controles establecidos en las normas \ISO/IEC 27001 e \ISO/IEC 27017. Como resultado, se estableció una base técnica y 
metodológica que permitió orientar la selección de tecnologías y la definición de aspectos de seguridad para la infraestructura cloud del 
grupo \GRID.

\textbf{Actividades realizadas}:
\begin{itemize}
   	\item Caracterización de la infraestructura tecnológica y organizacional del grupo \GRID.
	\item Identificación de actores, procesos y componentes relacionados con la seguridad de la infraestructura cloud.
	\item Aplicación de instrumentos de recolección de información orientados al análisis de seguridad de la infraestructura.
	\item Análisis de necesidades, problemas y oportunidades (NPO) relacionados con la ciberseguridad del entorno \GRID.
	\item Evaluación del nivel de alineación de la infraestructura con lineamientos de las normas \ISO/IEC 27001 e \ISO/IEC 27017.
\end{itemize}

\textbf{Resultados obtenidos}:
\begin{itemize}
        \item Identificación de necesidades relacionadas con la centralización y formalización de controles de seguridad.
	\item Detección de problemáticas asociadas a la baja adopción de estándares y buenas prácticas de seguridad de la información.
	\item Identificación de oportunidades para implementar mecanismos de Gestión Unificada de Amenazas (UTM).
	\item Reconocimiento de la necesidad de fortalecer el monitoreo y la gestión centralizada de eventos de seguridad.
	\item Definición de oportunidades de alineación con controles de las normas \ISO/IEC 27001 e \ISO/IEC 27017.
\end{itemize}

\subsection{Objetivo específico 2}
\noindent

\textbf{Enunciado}: Caracterizar referentes metodológicos y tecnológicos relacionados con la computación en la nube del \GRID.

\textbf{Estado de cumplimiento}: \textbf{CUMPLIDO COMPLETAMENTE}

El presente objetivo se cumplió mediante el desarrollo de un estudio de mapeo sistemático (SMS) orientado a identificar, analizar y 
clasificar referentes metodológicos y tecnológicos relacionados con arquitecturas de seguridad, monitoreo y protección de infraestructuras 
de computación en la nube. El proceso permitió consolidar una base de conocimiento enfocada en mecanismos, controles, tecnologías y 
enfoques utilizados en entornos académicos e institucionales similares al contexto del grupo \GRID.

A través del análisis de la literatura científica se identificaron tendencias relacionadas con arquitecturas de defensa en profundidad, 
sistemas SIEM, tecnologías IDPS, monitoreo continuo, correlación de eventos y enfoques de seguridad como Zero Trust. Asimismo, se 
evidenció la recurrencia de controles alineados con estándares internacionales como la \ISO/IEC 27001 y la \ISO/IEC 27017, proporcionando 
una base metodológica para la definición de aspectos de seguridad aplicables al entorno cloud del \GRID.

De igual manera, la caracterización de los referentes permitió construir una taxonomía orientada a clasificar enfoques de seguridad, tipos 
de contribución y contextos de aplicación identificados en la literatura, consolidando así un soporte técnico y conceptual para las fases 
posteriores de análisis DAR, selección tecnológica y modelado arquitectónico de la propuesta desarrollada.

\textbf{Actividades realizadas}:
\begin{itemize}
    \item Definición del protocolo metodológico para el desarrollo del estudio de mapeo sistemático (SMS).
	\item Construcción y aplicación de cadenas de búsqueda en bases de datos científicas relacionadas con seguridad y computación en la nube.
	\item Aplicación de criterios de inclusión, exclusión y evaluación de calidad de los estudios identificados.
	\item Análisis y clasificación de enfoques metodológicos, arquitecturas y tecnologías de seguridad presentes en la literatura científica.
	\item Construcción de una taxonomía orientada a categorizar los referentes metodológicos y tecnológicos identificados.
\end{itemize}

\textbf{Resultados obtenidos}:
\begin{itemize}
    \item Identificación de referentes metodológicos y tecnológicos aplicados a entornos de computación en la nube.
	\item Caracterización de arquitecturas de seguridad, controles y mecanismos de protección utilizados en infraestructuras cloud.
	\item Identificación de enfoques relacionados con monitoreo continuo, correlación de eventos, SIEM, IDPS y Zero Trust.
	\item Construcción de una taxonomía de enfoques de seguridad y contextos de aplicación en entornos académicos e institucionales.
	\item Consolidación de una base de conocimiento para apoyar la selección de tecnologías y la definición de aspectos de seguridad para el \GRID.
\end{itemize}

\subsection{Objetivo específico 3}
\noindent

\textbf{Enunciado}: Seleccionar referentes metodológicos y tecnológicos relacionados con la seguridad en la infraestructura de la 
computación en la nube del \GRID.

\textbf{Estado de cumplimiento}: \textbf{CUMPLIDO COMPLETAMENTE}

El presente objetivo se cumplió mediante la aplicación del proceso de análisis y resolución de decisiones (\DAR), el cual permitió evaluar 
y seleccionar referentes metodológicos y tecnológicos orientados al fortalecimiento de la seguridad dentro de la infraestructura de 
computación en la nube del \GRID. Para ello, se definieron criterios técnicos, operacionales y de alineación normativa relacionados con 
monitoreo, detección de amenazas, correlación de eventos, interoperabilidad, escalabilidad y capacidad de integración con entornos cloud.

A partir de este proceso, se analizaron diferentes alternativas tecnológicas asociadas a sistemas SIEM, herramientas IDPS y soluciones de 
seguridad perimetral, permitiendo identificar aquellas con mayor adecuación frente a las necesidades, problemas y oportunidades previamente 
detectadas en el entorno del \GRID. Como resultado, las tecnologías Suricata y Wazuh fueron seleccionadas como componentes principales para 
la detección de intrusiones, monitoreo continuo y correlación centralizada de eventos de seguridad.

Asimismo, el análisis permitió consolidar referentes metodológicos relacionados con arquitecturas de defensa en profundidad, monitoreo de 
seguridad y gestión centralizada de eventos, proporcionando soporte para la definición de la arquitectura propuesta y su alineación con 
controles de las normas \ISO/IEC 27001 e \ISO/IEC 27017.

\textbf{Actividades realizadas}:
\begin{itemize}
    \item Definición de criterios de evaluación para el análisis de tecnologías de seguridad.
	\item Aplicación de la metodología DAR para el análisis y resolución de decisiones tecnológicas.
	\item Evaluación comparativa de herramientas SIEM, IDPS y soluciones de seguridad perimetral.
	\item Análisis de capacidades de interoperabilidad, monitoreo y correlación de eventos de las tecnologías evaluadas.
	\item Revisión del alineamiento de las tecnologías con controles de la \ISO/IEC 27001 y la \ISO/IEC 27017.
\end{itemize}

\textbf{Resultados obtenidos}:
\begin{itemize}
    \item Selección de la tecnología Suricata como solución IDPS para monitoreo y detección de intrusiones.
	\item Selección de la plataforma Wazuh como solución SIEM para correlación y gestión centralizada de eventos.
	\item Identificación del SonicWall NSA 6600 como solución UTM para fortalecimiento del perímetro de seguridad.
	\item Consolidación de referentes metodológicos asociados a arquitecturas de defensa en profundidad y monitoreo continuo.
	\item Definición de una base tecnológica orientada al fortalecimiento de la seguridad de la infraestructura cloud del \GRID.
\end{itemize}

\subsection{Objetivo específico 4}
\noindent

\textbf{Enunciado}: Diseñar una arquitectura de seguridad de la infraestructura de la computación en la nube del grupo de investigación. 
\GRID.

\textbf{Estado de cumplimiento}: \textbf{CUMPLIDO COMPLETAMENTE}

El presente objetivo se cumplió mediante el diseño y modelado de una arquitectura de defensa en profundidad orientada al fortalecimiento de 
la seguridad de la infraestructura de computación en la nube del \GRID. La arquitectura propuesta integra mecanismos de monitoreo, 
correlación de eventos, detección de intrusiones, control perimetral y respuesta activa, permitiendo establecer un modelo multicapa 
enfocado en la supervisión continua y protección de los activos de información.

Para el desarrollo de la arquitectura se adoptó TOGAF como marco de referencia metodológico y Archimate como lenguaje de modelado 
empresarial, permitiendo representar de manera estructurada la interacción entre las capas de negocio, aplicación y tecnología. De manera 
complementaria, se utilizó Draw.io para el modelado detallado de la infraestructura, especificando zonas de red, dispositivos, segmentación 
lógica e interoperabilidad entre los componentes de seguridad implementados.

La arquitectura diseñada integra tecnologías como el firewall SonicWall NSA 6600, el IDPS Suricata, el SIEM Wazuh y mecanismos de 
segmentación y monitoreo continuo, estableciendo interoperabilidad entre componentes físicos y virtualizados. Asimismo, el modelo propuesto 
permite fortalecer capacidades de trazabilidad, gestión de incidentes, supervisión de eventos y protección perimetral dentro del entorno 
cloud del \GRID, manteniendo alineamiento con controles de las normas \ISO/IEC 27001 e \ISO/IEC 27017.

\textbf{Actividades realizadas}:
\begin{itemize}
    \item Definición de la arquitectura de defensa en profundidad para el entorno \GRID.
	\item Modelado arquitectónico mediante Archimate de las capas de negocio, aplicación y tecnología.
	\item Diagramación de la infraestructura y segmentación de red utilizando Draw.io.
	\item Integración conceptual de tecnologías SIEM, IDPS y UTM dentro de la arquitectura propuesta.
	\item Definición de flujos de monitoreo, correlación de eventos y respuesta activa.
	\item Relación de la arquitectura con controles de las normas \ISO/IEC 27001 e \ISO/IEC 27017.
\end{itemize}

\textbf{Resultados obtenidos}:
\begin{itemize}
	\item Diseño de una arquitectura de defensa en profundidad orientada a la protección de la infraestructura cloud del \GRID.
	\item Construcción de siete vistas arquitectónicas representadas mediante Archimate.
	\item Integración de mecanismos de monitoreo continuo, correlación de eventos y respuesta activa ante incidentes.
	\item Definición de interoperabilidad entre SonicWall NSA 6600, Suricata y Wazuh.
	\item Especificación de zonas de red, segmentación lógica y flujos de comunicación de la arquitectura.
	\item Consolidación de una arquitectura alineada con controles de seguridad de la \ISO/IEC 27001 y la \ISO/IEC 27017.
\end{itemize}

\subsection{Objetivo específico 5}
\noindent

\textbf{Enunciado}: Implementar un prototipo funcional del diseño de arquitectura de seguridad perteneciente a la infraestructura de la 
computación en la nube del \GRID.

\textbf{Estado de cumplimiento}: \textbf{CUMPLIDO COMPLETAMENTE}



\textbf{Actividades realizadas}:
\begin{itemize}
	\item Creación del entorno prototipo basado en la infraestructura tecnológica del \GRID.
	\item Despliegue de máquinas virtuales para la implementación de Suricata y Wazuh.
	\item Configuración del firewall SonicWall NSA 6600 para control y monitoreo del tráfico de red.
	\item Implementación de mecanismos de Port Mirroring y segmentación de red mediante el switch Ruijie.
	\item Integración funcional entre Suricata, Wazuh y SonicWall NSA 6600.
	\item Configuración de mecanismos de monitoreo continuo, correlación de eventos y respuesta activa.
\end{itemize}

\textbf{Resultados obtenidos}:
\begin{itemize}
	\item Implementación funcional de un prototipo de arquitectura de defensa en profundidad para el entorno \GRID.
	\item Validación operativa de procesos de monitoreo y detección de eventos de seguridad.
	\item Centralización y correlación de registros mediante la plataforma SIEM Wazuh.
	\item Verificación de capacidades de inspección y análisis de tráfico utilizando Suricata.
	\item Integración del SonicWall NSA 6600 como mecanismo de protección perimetral y control de acceso.
	\item Comprobación de la interoperabilidad entre componentes físicos y virtualizados dentro del entorno prototipo.
\end{itemize}

\subsection{Objetivo específico 6}
\noindent

\textbf{Enunciado}: Validar la arquitectura de seguridad mediante la implementación del prototipo de arquitectura funcional en relación 
con la computación en la nube del grupo \GRID.

\textbf{Estado de cumplimiento}: \textbf{CUMPLIDO COMPLETAMENTE}



\textbf{Actividades realizadas}:
\begin{itemize}
	\item Integración del SonicWall NSA 6600 como mecanismo de protección perimetral y control de acceso.
	\item Comprobación de la interoperabilidad entre componentes físicos y virtualizados dentro del entorno prototipo.
\end{itemize}

\textbf{Resultados obtenidos}:
\begin{itemize}
	\item Integración del SonicWall NSA 6600 como mecanismo de protección perimetral y control de acceso.
	\item Comprobación de la interoperabilidad entre componentes físicos y virtualizados dentro del entorno prototipo.
\end{itemize}

\section{Análisis del grado de cumplimiento}
\noindent

\subsection{Cumplimiento cuantitativo}
\noindent

El análisis cuantitativo del cumplimiento de objetivos muestra:

\begin{itemize}
    \item \textbf{Objetivo general}: 100\% cumplido.
    \item \textbf{Objetivos específicos}: 100\% de cumplimiento en los 5 objetivos planteados.
    \item \textbf{Entregables}: 100\% de entregables completados según cronograma establecido.
    \item \textbf{Casos de prueba}: 100\% de casos de prueba ejecutados exitosamente.
\end{itemize}

\subsection{Cumplimiento cualitativo}
\noindent

Desde una perspectiva cualitativa, el proyecto permitió consolidar una propuesta orientada al fortalecimiento de la seguridad de la 
infraestructura de computación en la nube del \GRID, evidenciando resultados relevantes tanto a nivel metodológico como tecnológico:

\textbf{Integralidad de la solución}: Los aspectos de seguridad propuestos incorporan mecanismos de protección perimetral, monitoreo 
continuo, correlación de eventos y respuesta activa, permitiendo abordar la seguridad desde un enfoque de defensa en profundidad.

\textbf{Alineamiento normativo}: El diseño arquitectónico y las tecnologías seleccionadas mantienen relación con controles de las normas 
\ISO/IEC 27001 e \ISO/IEC 27017, fortaleciendo la gestión de seguridad y la formalización de controles dentro del entorno cloud del \GRID.

\textbf{Interoperabilidad tecnológica}: La integración entre SonicWall NSA 6600, Suricata, Wazuh y los mecanismos de segmentación de red 
permitió establecer capacidades coordinadas de supervisión, detección y análisis de eventos de seguridad.

\textbf{Modelado y trazabilidad}: El uso de Archimate y Draw.io permitió representar de manera estructurada las relaciones entre capas de 
negocio, aplicación y tecnología, facilitando la comprensión operacional de la arquitectura propuesta.

\textbf{Reproducibilidad}: La documentación generada, junto con las vistas arquitectónicas y especificaciones técnicas desarrolladas, 
proporciona una base metodológica y técnica que facilita la futura implementación, escalabilidad y adaptación de la arquitectura en 
contextos similares.

\section{Impacto en las NPO identificadas}
\noindent

\subsection{Necesidades atendidas}
\noindent

\textbf{Centralización de controles de seguridad}: Los aspectos de seguridad propuesta integran mecanismos de monitoreo, correlación de 
eventos, control perimetral y respuesta activa dentro de un ecosistema tecnológico unificado, permitiendo fortalecer la gestión 
centralizada de la seguridad en la infraestructura cloud del \GRID.

\textbf{Fortalecimiento del monitoreo y la supervisión}: La integración de tecnologías como Wazuh y Suricata permite establecer capacidades 
de monitoreo continuo, análisis de tráfico y correlación de eventos orientadas a la detección de comportamientos anómalos dentro de la 
infraestructura.

\textbf{Protección perimetral y segmentación}: La incorporación del SonicWall NSA 6600 y mecanismos de segmentación lógica contribuye al 
fortalecimiento del perímetro de seguridad y al control de acceso entre las diferentes zonas de red del entorno \GRID.

\textbf{Alineamiento con estándares internacionales}: El diseño de la arquitectura y la selección de tecnologías se orientaron al 
fortalecimiento de controles relacionados con las normas \ISO/IEC 27001 e \ISO/IEC 27017, promoviendo mejores prácticas de seguridad de la 
información.

\textbf{Formalización y trazabilidad}: El modelado arquitectónico realizado mediante Archimate y Draw.io permitió documentar de forma 
estructurada la interacción entre procesos, tecnologías y actores involucrados en la operación de la infraestructura de seguridad propuesta.

\subsection{Oportunidades aprovechadas}
\noindent

\textbf{Implementación de una arquitectura de defensa en profundidad}: El proyecto permitió diseñar un prototipo de una arquitectura de 
defensa en profundidadorientada al fortalecimiento de la protección, monitoreo y supervisión de la infraestructura cloud del \GRID\ 
mediante mecanismos multicapa de seguridad.

\textbf{Integración de tecnologías de seguridad}: La incorporación de herramientas como SonicWall NSA 6600, Suricata y Wazuh permitió 
consolidar capacidades de correlación de eventos, monitoreo continuo y respuesta activa dentro de un entorno centralizado.

\textbf{Fortalecimiento de la gobernanza de seguridad}: La alineación del prototipo de la arquitectura en seguridad con controles de las 
normas \ISO/IEC 27001 e \ISO/IEC 27017 permitió establecer una base orientada a la formalización y mejora continua de la seguridad de la 
información en el \GRID.

\textbf{Capacidades de trazabilidad y supervisión}: El prototipo de la arquitectura propuesta habilita una mayor visibilidad sobre los 
eventos y actividades ejecutadas dentro de la infraestructura, fortaleciendo la capacidad de análisis y gestión de incidentes de seguridad.

\textbf{Base para futuras implementaciones}: El modelado arquitectónico, la documentación técnica y el entorno prototipo desarrollado 
proporcionan una base reproducible y escalable para futuras implementaciones de seguridad en contextos académicos similares.

\subsection{Problemas resueltos}
\noindent

\textbf{Baja formalización de mecanismos de seguridad}: El modelado arquitectónico y la alineación con controles de la \ISO/IEC 27001 y la
\ISO/IEC 27017 permitieron estructurar de forma más organizada los procesos y controles relacionados con la seguridad del entorno \GRID.

\textbf{Limitaciones en monitoreo y trazabilidad}: La integración de Wazuh y Suricata permitió fortalecer la capacidad de supervisión 
continua, análisis de tráfico y detección de eventos sospechosos dentro de la infraestructura.

\textbf{Falta de centralización de eventos}: La implementación de un SIEM permitió consolidar registros y eventos provenientes de 
múltiples fuentes dentro de una plataforma centralizada de monitoreo.

\textbf{Ausencia de una arquitectura de seguridad estructurada}: El diseño de la arquitectura de defensa en profundidad permitió 
establecer un modelo organizado de interacción entre capas de negocio, aplicación y tecnología para la protección de la infraestructura 
cloud del \GRID.

\section{Lecciones aprendidas}
\noindent

\subsection{Aspectos metodológicos}
\noindent

La integración del análisis de necesidades, problemas y oportunidades (\NPO), el estudio de mapeo sistemático (\SMS), el análisis \DAR\ y 
el modelado arquitectónico permitió estructurar una metodología orientada a la toma de decisiones fundamentadas para el fortalecimiento de 
la seguridad del entorno \GRID. Asimismo, se evidenció la importancia de combinar enfoques cualitativos, normativos y tecnológicos para 
construir propuestas alineadas tanto con estándares internacionales como con las necesidades reales de la infraestructura.

\subsection{Aspectos técnicos}
\noindent

La interoperabilidad entre tecnologías como SonicWall NSA 6600, Suricata y Wazuh demostró la viabilidad de implementar capacidades de 
monitoreo, correlación de eventos y respuesta activa dentro de un entorno basado en componentes físicos y virtualizados. De igual manera, 
el modelado arquitectónico mediante Archimate y Draw.io facilitó la comprensión estructurada de la interacción entre capas de negocio, 
aplicación y tecnología, fortaleciendo la trazabilidad operacional de la arquitectura propuesta.

\subsection{Aspectos institucionales}
\noindent

El desarrollo del proyecto permitió evidenciar la importancia de alinear los mecanismos de seguridad con las necesidades operativas y 
académicas del \GRID, especialmente en aspectos relacionados con monitoreo, protección de activos y supervisión continua de la 
infraestructura cloud. Asimismo, se identificó que la formalización de controles y la adopción de estándares como la \ISO/IEC 27001 y la 
\ISO/IEC 27017 constituyen elementos clave para fortalecer la madurez de seguridad de la información dentro de entornos académicos y de 
investigación.

\section{Conclusión del cumplimiento de objetivos}
\noindent

La evaluación realizada permite concluir que los objetivos planteados para la presente investigación fueron cumplidos satisfactoriamente, 
permitiendo identificar necesidades, problemas y oportunidades relacionadas con la seguridad de la infraestructura cloud del \GRID, así 
como seleccionar tecnologías y enfoques metodológicos acordes con el contexto institucional y académico del grupo de investigación.

De igual manera, el desarrollo de la arquitectura de defensa en profundidad y su representación mediante modelos arquitectónicos permitió 
consolidar una propuesta estructurada, alineada con controles de las normas \ISO/IEC 27001 e \ISO/IEC 27017 y orientada al fortalecimiento 
de la protección, trazabilidad y gestión de la seguridad dentro del entorno \GRID.

Finalmente, los resultados obtenidos proporcionan una base técnica, metodológica y documental que facilita futuras implementaciones y 
procesos de mejora continua, contribuyendo al fortalecimiento de la seguridad de la información en entornos académicos y de investigación 
con características similares.
